\section{Lecture 21: Six Functors for Analytic Stacks}
\label{lec:21}

The bulk of this lecture sets up the six-functor formalism for analytic stacks
and, following Liu--Zheng and Mann, shows that the formalism itself \emph{dictates}
the Grothendieck topology: there is a smallest class of ``shriekable'' morphisms
of analytic stacks to which the six operations extend. Along the way Scholze
gives inductive definitions of \emph{proper} and \emph{cohomologically smooth}
maps and computes $f_!$ in examples (an interval, the real line). The lecture
opens by discharging a promise from \cref{lec:19,lec:20}---the explicit
construction of the norm on the gaseous base ring---and closes by returning to
the Tate elliptic curve $E_q$, now shown to be proper and smooth over the
gaseous base and then \emph{algebraised} into an honest scheme.

Throughout, ``ring'' means commutative ring and everything is (light) condensed.

\subsection{Construction of the norm}
\label{subsec:21:norm}

We first construct the norm on the gaseous base ring
\begin{equation}
\label{eq:21:gaseous-base}
A\ =\ \mathbb Z[\widehat q^{\pm1}]_{\gasx},\qquad |q|=\tfrac12,
\end{equation}
promised by Clausen in \cref{lec:20} (existence half of
\cref{thm:20:classification}). By \cref{prop:20:gaseous-factorization} one may
assume all rational powers of $q$ are available. The task is to produce the norm
$N\colon\mathbb P^1_A\to[0,\infty]$ of \cref{def:20:norm}; by
\cref{rmk:20:overconvergent-disks} this amounts to specifying, compatibly, the
idempotent algebras cutting out the loci $\{|T|\le r\}$, i.e.\ the
\emph{overconvergent} functions on the closed disc of radius $r$.

As Clausen explained, the norm axioms force the answer up to a single natural
guess. The naive candidate is the free algebra on a topologically nilpotent
generator obtained by rescaling the coordinate $T$ by a power of $q$; this is
not quite right, but becomes correct after passing to the overconvergent
(colimit) version.

\begin{definition}[Overconvergent closed disc]
\label{def:21:overconvergent-disc}
For $r\in\mathbb R_{>0}$ set
\begin{equation}
\label{eq:21:disc-algebra}
\mathcal O\bigl(\{|T|\le r\}\bigr)^{t}\ :=\
\varinjlim_{t\,:\,2^{t}>r}\ A\bigl[\widehat{q^{t}T}\bigr],
\end{equation}
where $A[\widehat{s}]$ denotes the free $A$-algebra on a topologically nilpotent
element $s$ (the ring $P_A$ of \cref{rmk:7:P-ring}, with $s$ acting as the shift),
and the colimit runs over rational $t$. The superscript $t$ records that this is
the \emph{overconvergent} structure sheaf, as in \cref{eq:19:overconvergent}.
\end{definition}

The scaling is dictated by wanting $|T|=r$ and $|q^{t}T|<1$ (so that $q^{t}T$ is
genuinely topologically nilpotent): since $|q|=\tfrac12$, the condition
$(\tfrac12)^{t}\,r=|q^{t}T|<1$ reads $2^{t}>r$. This is the only possible choice.
One must then check that these algebras are idempotent and that (together with
the analogous algebras ``at $\infty$'', obtained by replacing $T$ by $T^{-1}$)
they assemble into a map $\mathbb P^1_A\to[0,\infty]$; all of these are
straightforward computations. We record the one that matters.

\begin{proposition}[Discs of different radii are disjoint]
\label{prop:21:disjoint-discs}
For $r<s$ the closed disc $\{|T|\le r\}$ and the complement of the open disc
$\{|T|\ge s\}$ do not intersect: with the algebra
$\mathcal O(\{|T|\ge s\})^{t}=\varinjlim A[T,\widehat{q^{s'}T^{-1}}]$ obtained by
inversion,
\begin{equation}
\label{eq:21:disjoint-discs}
A\bigl[\widehat{q^{t}T},\,T^{-1}\bigr]\ \otimes_{A[T^{\pm1}]}\
A\bigl[T,\,\widehat{q^{s}T^{-1}}\bigr]\ =\ 0
\qquad(t<0,\ s<0).
\end{equation}
\end{proposition}

\begin{proof}[Sketch]
Writing $x=q^{t}T$ and $y=q^{s}T^{-1}$, one has $xy=q^{t+s}$, so the tensor
product is the free algebra on the two topologically nilpotent generators
$x,y$ subject to the single relation coming from their product,
\begin{equation}
\label{eq:21:idempotent-collapse}
A[\widehat x,\widehat y]\big/\bigl(1-q^{a}xy\bigr)\ =\ P_A\big/\bigl(1-q^{a}\,\mathrm{shift}\bigr),
\qquad a=-(t+s)>0 .
\end{equation}
But the gaseous ring structure was \emph{defined} by declaring
$1-q\cdot\mathrm{shift}$ (and hence, since the exponent is immaterial,
$1-q^{a}\cdot\mathrm{shift}$) to be an isomorphism on the compact projective
generator $P$ (\cref{thm:14:gaseous-base}). Thus the quotient vanishes.
\end{proof}

\begin{remark}[Two honest-subspace variants of the disc]
\label{rmk:21:disc-variants}
The basic disc used above,
\begin{equation}
\label{eq:21:basic-disc}
D_A\ =\ \AnSpec\bigl(A[\widehat T]\bigr)\ \longrightarrow\ \mathbb P^1_A ,
\end{equation}
is proper but not a monomorphism (\cref{def:20:analytic-affine-line}). Scholze
noted a cleaner perspective: instead of forcing $T$ topologically nilpotent one
can force it merely \emph{gaseous}, i.e.\ pass to the analytic ring
$A[T]^{T\text{-}\gasx}$ obtained by enforcing $1-T\cdot\mathrm{shift}$ to be an
isomorphism on $P$ base-changed to it. Then $\AnSpec(A[T]^{T\text{-}\gasx})$
\emph{is} an honest subspace of $\mathbb P^1_A$, sitting between the open and
closed unit discs,
\begin{equation}
\label{eq:21:sandwich}
N^{-1}\bigl([0,1)\bigr)\ \subseteq\ \AnSpec\bigl(A[T]^{T\text{-}\gasx}\bigr)\
\subseteq\ N^{-1}\bigl([0,1]\bigr).
\end{equation}
One may also take the intersection of $D_A$ with this gaseous disc, or the
variant $A[\widehat T, T]^{T\text{-}\gasx}$; all three are again sandwiched
between the overconvergent closed disc and the open disc, so for the
\emph{overconvergent} construction the choice is immaterial. This gives a
calculation-free route to idempotency: two of the candidates are sandwiched in
each other within the colimit, and since one term is always a monomorphism the
idempotency for the other follows formally. The one genuinely necessary
computation is \cref{prop:21:disjoint-discs}, that the closed disc
$\{|T|\le r\}$ around $0$ and the disc $\{|T|\ge s\}$ around $\infty$ do not
meet when $r<s$.
\end{remark}

\subsection{Six functors: recollection and correspondences}
\label{subsec:21:six-functors}

Recall the two functors on analytic rings assembled so far. The primary
invariant is
\begin{equation}
\label{eq:21:two-functors}
A\ \longmapsto\ D(A),\qquad A\ \longmapsto\ \PrL_{D(A)}=\Mod_{D(A)}(\PrL),
\end{equation}
a functor $\AnRing\to\CAlg(\Catinf)$ landing in symmetric monoidal
$\infty$-categories: the algebra structure supplies the tensor product $\otimes$,
which is closed, so there is an internal $\RHom$; and functoriality in the ring
gives, for each map of rings, a base-change (pullback) functor $f^{*}$ with right
adjoint $f_{*}$. These are four of the six operations. On the class of
$!$-able (shriekable) maps (\cref{def:16:shriekable,def:18:shriekable}) we
constructed in addition the lower-shriek $f_{!}$ and its right adjoint $f^{!}$.

A convenient way to package all this is the notion of a six-functor formalism
via correspondence categories.

\begin{definition}[Correspondence category]
\label{def:21:correspondence}
Let $\mathcal C$ be an $\infty$-category with all finite limits and let
$E$ be a class of morphisms stable under composition (including identities) and
base change. The \emph{correspondence category} $\Corr(\mathcal C,E)$ is the
symmetric monoidal $\infty$-category with the same objects as $\mathcal C$,
whose morphisms $X\to Y$ are spans
\begin{equation}
\label{eq:21:span}
\begin{tikzcd}[column sep=small, row sep=small]
& W \arrow[dl] \arrow[dr, "\in E"] & \\
X & & Y
\end{tikzcd}
\end{equation}
with the right leg in $E$, composed by forming the fibre product
\begin{equation}
\label{eq:21:span-composition}
\begin{tikzcd}[column sep=small, row sep=small]
& & W\times_Y W' \arrow[dl] \arrow[dr, "\in E"] & & \\
& W \arrow[dl] \arrow[dr, "\in E"] & & W' \arrow[dl] \arrow[dr, "\in E"] & \\
X & & Y & & Z
\end{tikzcd}
\end{equation}
(the outer legs form the composite span), and with symmetric monoidal structure
$X\otimes Y=X\times Y$ given by the product in $\mathcal C$.
\end{definition}

\begin{definition}[Three- and six-functor formalisms]
\label{def:21:functor-formalism}
A \emph{$3$-functor formalism} on $(\mathcal C,E)$---encoding $\otimes$, $f^{*}$
and $f_{!}$ (for $f\in E$)---is a lax symmetric monoidal functor
\begin{equation}
\label{eq:21:3-functor}
D\colon\ \bigl(\Corr(\mathcal C,E),\otimes\bigr)\ \longrightarrow\
\bigl(\Catinf,\times\bigr),
\end{equation}
the lax structure being the natural maps
$D(X)\times D(Y)\to D(X\times Y)$ (in one direction only, not necessarily
isomorphisms). Applied to the empty product, the lax structure endows
$D(\ast)=D(\text{final object})$ with a distinguished (unit) object, and every
$D(X)$ receives a unit by pullback from $\ast$. It is a \emph{$6$-functor
formalism} if $\RHom$, $f_{*}$ and $f^{!}$ (for $f\in E$) exist as the
corresponding right adjoints.
\end{definition}

Producing such a functor directly is a formidable amount of coherence data.
Conveniently, Liu--Zheng gave an extremely simple combinatorial recipe (the
difficulty being entirely in the proof that it works), which makes these objects
easy to construct in practice. For $\mathcal C=(\AnRing)^{\mathrm{op}}$ and $E$
the shriekable maps, $A\mapsto D(A)$ (equivalently $A\mapsto\PrL_{D(A)}$) is such
an example.

\subsection{The topology dictated by the formalism}
\label{subsec:21:topology}

We now extend the formalism from affine analytic spaces to all analytic stacks,
enlarging the class of shriekable maps as little as possible.

\begin{theorem}[Liu--Zheng, Mann]
\label{thm:21:minimal-class}
There is a (unique) smallest class $\widetilde E$ of morphisms in $\AnStk$ with
the following properties.
\begin{enumerate}
\item $E\subseteq\widetilde E$, where $E$ is the class of shriekable maps of
analytic rings.
\item \textup{(Stable under disjoint unions.)} If $f_i\colon X_i\to Y$ lie in
$\widetilde E$, then so does $f=\coprod_i f_i\colon\coprod_i X_i\to Y$.
\item \textup{($\widetilde E$ local on the target.)} If $f\colon X\to Y$ in
$\AnStk$ admits a cover $\coprod_i Y_i'\twoheadrightarrow Y$ such that each base
change $f_i'=f\times_Y Y_i'\colon X\times_Y Y_i'\to Y_i'$ lies in $\widetilde E$,
then $f\in\widetilde E$.
\item \textup{($\widetilde E$ local on the source.)} If $f\colon X\to Y$ admits a
cover $g\colon\widetilde X\twoheadrightarrow X$ with $g\in\widetilde E$ of
universal $!$-descent and with the composite $\widetilde X\to Y$ in
$\widetilde E$, then $f\in\widetilde E$.
\item If $f\colon X\to Y\in\widetilde E$ with $Y$ affine, then there is a cover
$g\colon\coprod_i X_i\twoheadrightarrow X$ with $g\in\widetilde E$, all $X_i$
affine, and each composite $X_i\to Y$ in $E$.
\item The $6$-functor formalism extends uniquely from
$\bigl((\AnRing)^{\mathrm{op}},E\bigr)$ to $(\AnStk,\widetilde E)$.
\end{enumerate}
\end{theorem}

The theorem follows from results of Liu--Zheng \cite{liu2024enhancedoperationsbasechange}
and Mann \cite{mann2022padic6functorformalismrigidanalytic}; this is a general
fact about six-functor formalisms, independent of the specifics of the present
one, and an account is given in Scholze's notes on six functors
\cite{scholze2023six} in a slightly different setup.

\begin{remark}[Localising on source and target]
\label{rmk:21:localise}
Property~(3) lets one check shriekability after localising on the target, and
property~(4) after localising (by a shriekable cover) on the source; iterating
the two---each new shriekable cover produces further maps to be covered---one
reduces to affine situations. Consequently, checking that a given map is
shriekable is in practice easy: essentially every morphism is, and the local
criteria make it routine to verify. The lower-shriek on $Y$ is then reconstructed
from the one on a cover $Y'$: since $f_!$ commutes with base change, defining it
after base change to $Y'$ automatically produces the descent datum along all
further base changes.
\end{remark}

\begin{example}[Reconstructing $f_!$ from a cover]
\label{ex:21:reconstruct-fshriek}
Suppose $X$ is stacky, presented by a cover $\widetilde X\twoheadrightarrow X$
over $Y$ for which shriek functors are already available. Then
\begin{equation}
\label{eq:21:D-limit-colimit}
D(X)\ =\ \varprojlim_{\Delta}{}^{!}\Bigl[\,D(\widetilde X)\rightrightarrows
D(\widetilde X\times_X\widetilde X)\Rrightarrow\cdots\Bigr]\ =\
\varinjlim_{\Delta^{\mathrm{op}}}{}^{!}\Bigl[\,D(\widetilde X)\rightrightarrows
D(\widetilde X\times_X\widetilde X)\Rrightarrow\cdots\Bigr]\quad\text{in }\PrL,
\end{equation}
the first a limit along the upper-shriek transition functors and the second the
same diagram read as a colimit along the lower-shriek functors in $\PrL$
(\cref{cor:17:descent-codescent}). As all the $D(\widetilde X^{\times_X n})$
carry lower-shriek functors to $D(Y)$ compatible with the transition maps, and
$D(Y)$ is presentable, the colimit yields a unique $f_!\colon D(X)\to D(Y)$.
\end{example}

\begin{example}[The interval, and compactly supported cohomology]
\label{ex:21:interval}
Consider the interval mapping to a point, $[0,1]\to\ast$, which lives in
condensed sets and hence (\cref{ex:19:condensed-anima}) in analytic stacks. By
Clausen's descendable cover (\cref{lec:20}) there is a Cantor set surjecting onto
$[0,1]$; base-changing $[0,1]\to\ast$ along it gives a light profinite set over
the Cantor set, which is a basic flat cover (a map of light profinite sets, in
$E$). Hence
\begin{equation}
\label{eq:21:interval-cover}
\{0,1\}^{\mathbb N}\ \xrightarrow{\ \in\widetilde E\ }\ [0,1]\
\xrightarrow{\ f\ }\ \ast,
\end{equation}
with the composite $\{0,1\}^{\mathbb N}\to\ast$ (a map of light profinite sets)
in $E$, exhibits $f$ as shriekable. Although we only
defined shriek functors for maps of light profinite sets, the formalism now
produces $f_!$ for the interval, and one observes
\begin{equation}
\label{eq:21:interval-fshriek}
f_!\ =\ f_* ,
\end{equation}
as one expects since $[0,1]\to\ast$ is ``proper''; this $f_!$ is compactly
supported cohomology of $[0,1]$.
\end{example}

\begin{example}[The real line]
\label{ex:21:real-line}
Locally compact Hausdorff spaces are handled the same way. The real line is
covered by the disjoint union of the closed intervals,
\begin{equation}
\label{eq:21:real-cover}
\coprod_{n}[-n,n]\ \xrightarrow{\ \in\widetilde E\ }\ \mathbb R ,
\end{equation}
which is surjective (as light condensed sets, $\mathbb R$ is this union) and lies
in $\widetilde E$: covering by disjoint unions of profinite sets exhibits it as a
disjoint union of closed subsets, each already in $\widetilde E$. The formalism
thus yields an $f_!$ for $\mathbb R\to\ast$, and one checks it is the usual
lower-shriek. Indeed $D(\mathbb R)$ is the colimit
$\varinjlim_n D([-n,n])$ along the lower-shriek maps, so it suffices to describe
$f_!$ on each $[-n,n]$, where it is compactly supported cohomology. This abstract
procedure recovers a great deal of the geometric intuition about lower-shriek.
\end{example}

\subsection{Proper and smooth maps}
\label{subsec:21:proper-smooth}

We isolate two important classes: proper maps (for which $f_!=f_*$) and
cohomologically smooth maps (for which $f^!$ is a twist of $f^*$).

\begin{definition}[Proper map]
\label{def:21:proper}
A morphism $f\colon X\to Y$ in $\widetilde E$, assumed quasicompact and
quasiseparated, is \emph{proper} if one of the following (inductive) conditions
holds.
\begin{enumerate}
\item \textup{(Relatively representable in affines and proper.)} There is a cover
$\coprod_i Y_i'\twoheadrightarrow Y$ by affines such that each
$X_i'=X\times_Y Y_i'$ is affine and $X_i'\to Y_i'$ is proper in the earlier sense
of \cref{def:16:proper}, i.e.\ $X_i'$ carries the induced analytic ring
structure. In this case one obtains a canonical equivalence $f_!\cong f_*$.
\item \textup{(Proper diagonal.)} The diagonal
$\Delta_f\colon X\to X\times_Y X$ is proper (in this same sense, inductively).
\end{enumerate}
\end{definition}

\begin{remark}[The comparison map $f_!\to f_*$]
\label{rmk:21:comparison-map}
Under condition~(1) the identification $f_!\cong f_*$ is definitional. In
general it is built from a proper diagonal. Consider
\begin{equation}
\label{eq:21:diagonal-square}
\begin{tikzcd}[row sep=large, column sep=large]
X \arrow[r, "\Delta_f"] \arrow[dr, "f"'] & X\times_Y X \arrow[r, "p_2"] \arrow[d, "p_1"] & X \arrow[d, "f"] \\
& X \arrow[r, "f"'] & Y .
\end{tikzcd}
\end{equation}
Base change on the right-hand square gives $f^{*}f_{!}=p_{2!}p_1^{*}$, and,
inserting the unit $\mathbf 1\to\Delta_{f*}\Delta_f^{*}$ of the adjunction
$\Delta_f^{*}\dashv\Delta_{f*}$ (and using $\Delta_{f!}=\Delta_{f*}$, available
since $\Delta_f$ is proper),
\begin{equation}
\label{eq:21:comparison}
f^{*}f_{!}\ =\ p_{2!}p_1^{*}\ \longrightarrow\
p_{2!}\Delta_{f*}\Delta_f^{*}p_1^{*}\ \cong\
p_{2!}\Delta_{f!}\Delta_f^{*}p_1^{*}\ =\ \mathrm{id},
\end{equation}
whose adjoint is a natural transformation $f_{!}\to f_{*}$. One \emph{requires}
this to be an isomorphism; by a diagram chase (in the six-functor notes
\cite{scholze2023six}) it suffices to check it on the unit,
$f_!(\mathbf 1)\xrightarrow{\sim}f_*(\mathbf 1)$, and then it holds after every
base change. The class of proper maps is stable under base change and
composition. Scholze noted the definition should be restricted to quasicompact,
quasiseparated maps: all maps in the two conditions are automatically QCQS
(the basic ones are, and diagonals are), whereas passing to arbitrary disjoint
unions $\coprod_i(X_i\to Y_i)$ with the senses tending to infinity would fail.
\end{remark}

\begin{remark}[$f^!$ of the unit; a caveat for schemes]
\label{rmk:21:fshriek-unit}
Asked for an explicit description of $f^!(\mathbf 1)$ for $f\colon X\to\ast$: in
the examples above (interval, real line) $f_!$ is the usual lower-shriek, so
$f^!$ is the usual dualising complex, and in essentially any situation where a
dualising complex already exists, $f^!$ agrees with it. One caveat: if a scheme
is embedded into analytic stacks with its \emph{trivial} analytic ring structure
(\cref{ex:19:derived-schemes}), then every map is proper, $f_!=f_*$, and $f^!$ is
the poorly behaved right adjoint of $f_*$---not the Grothendieck-duality shriek.
To recover the usual dualising complex of a finite-type scheme one should instead
use the relatively-solid embedding (\cref{rmk:19:two-embeddings}).
\end{remark}

\begin{definition}[Cohomologically smooth map]
\label{def:21:smooth}
A morphism $f\colon X\to Y$ in $\widetilde E$ is \emph{cohomologically smooth} if,
after any base change, the natural map
\begin{equation}
\label{eq:21:smooth}
f^{*}(-)\otimes f^{!}(\mathbf 1)\ \xrightarrow{\ \sim\ }\ f^{!}(-)
\end{equation}
is an isomorphism, $f^{!}(\mathbf 1)\in D(X)$ is $\otimes$-invertible, and its
formation commutes with base change.
\end{definition}

The map \eqref{eq:21:smooth} is adjoint to the projection-formula counit
$f_!\bigl(f^{*}(-)\otimes f^{!}(\mathbf 1)\bigr)=(-)\otimes f_!f^{!}(\mathbf 1)
\to(-)$. Unlike properness, smoothness is not stable under passing to diagonals,
and one does not expect $f^{!}=f^{*}$ on the nose---only up to the invertible
twist $f^{!}(\mathbf 1)$, the relative dualising object. Cohomologically smooth
maps are stable under base change and composition, and cohomological smoothness
is local on the source (a cohomologically smooth cover of a cohomologically
smooth map descends).

\subsection{The Tate elliptic curve, revisited}
\label{subsec:21:tate}

Return to the gaseous base $A=\mathbb Z[\widehat q^{\pm1}]_{\gasx}$, now equipped
with the norm of \cref{subsec:21:norm}. The analytic multiplicative group is the
preimage of the open interval,
\begin{equation}
\label{eq:21:Gm-an}
{\Gm}^{\mathrm{an}}_A\ =\ N^{-1}\bigl((0,\infty)\bigr)\ \subseteq\ \mathbb P^1_A ,
\end{equation}
and the open inclusion $j\colon{\Gm}^{\mathrm{an}}_A\hookrightarrow\mathbb P^1_A$ is
shriekable; it is in fact an open immersion, with $j^{!}=j^{*}$ and
$j^{*}j_{!}=\mathrm{id}$, and it is cohomologically smooth. (Everything here
reduces, by localising at $0$ and $\infty$, to the usual six operations on
locally compact Hausdorff spaces.) The projection
$\pi\colon\mathbb P^1_A\to\AnSpec(A)$ is proper: it is not relatively
representable in affines, but its diagonal is (it is the diagonal of $T_1$ inside
$\mathbb P^1$, an affine situation), and one checks $\pi_!$ and $\pi_*$ of the
unit agree, exactly as in algebraic geometry. Moreover $\pi$ is cohomologically
smooth---since a morphism of six-functor formalisms (here schemes into analytic
stacks) preserves cohomologically smooth maps, the algebraic smoothness of
$\mathbb P^1$ implies the analytic one.

The action of $q$ by $T\mapsto qT$ is free and totally discontinuous, and one
forms the quotient in analytic stacks. Rescaling so that $|q|=\tfrac12$, the norm
descends to $(0,\infty)/(\tfrac12)^{\mathbb Z}\cong S^1$, giving
\begin{equation}
\label{eq:21:tate-tower}
\begin{tikzcd}[column sep=large]
E_{q,A}={\Gm}^{\mathrm{an}}_A/q^{\mathbb Z}
\arrow[r, "\text{proper}"] &
(0,\infty)/(\tfrac12)^{\mathbb Z}=S^1
\arrow[r, "\text{proper}"] &
\AnSpec(A).
\end{tikzcd}
\end{equation}
The first map is proper (checked by base change over $(0,\infty)$) and the second
is proper, so $E_{q,A}\to\AnSpec(A)$ is proper. The quotient map
${\Gm}^{\mathrm{an}}_A\to E_{q,A}$ is locally split, so $E_{q,A}$ is locally
isomorphic to ${\Gm}^{\mathrm{an}}_A$ and hence cohomologically smooth. Thus
$E_{q,A}$ is a proper smooth curve, the \emph{Tate elliptic curve}, carrying the
unit section $1\in\Gm^{\mathrm{an}}$.

\begin{remark}[Open immersions are not affine-representable]
\label{rmk:21:open-not-representable}
Asked whether open immersions of analytic stacks are relatively representable in
affines: they are not. Were $j$ representable in affines it would in particular be
quasicompact, whereas an open immersion such as
${\Gm}^{\mathrm{an}}_A\hookrightarrow\mathbb P^1_A$ need not be. The open
immersions of the present formalism are far more general than the classical ones.
\end{remark}

\subsection{Algebraisation}
\label{subsec:21:algebraisation}

It remains to see that the Tate curve is an honest (algebraic) elliptic curve.

\begin{theorem}[$E_q$ is algebraic]
\label{thm:21:tate-algebraic}
The Tate elliptic curve $E_{q,A}\in\AnStk_A$ lies in the essential image of the
fully faithful functor $\mathrm{Sch}_{A^{\triangleright}(\ast)}\to\AnStk_A$ from
schemes over the underlying discrete ring; i.e.\ it is the analytification of an
elliptic curve over $A^{\triangleright}(\ast)$.
\end{theorem}

This is an instance of a general algebraisation principle, whose proof is the
classical construction of $\operatorname{Proj}$ of a section ring, transported to
the analytic-stack setting.

\begin{proposition}[General algebraisation setup]
\label{prop:21:algebraisation}
Let $A\in\AnRing$ and let $X\in\AnStk_{/A}$ be proper. Suppose $X$ carries a line
bundle $L$ such that, after replacing $L$ by a tensor power, one has:
\begin{enumerate}
\item $L$ is globally generated---there is a surjection
$\mathcal O_X^{\oplus N}\twoheadrightarrow L$ (equivalently, a surjection on
$H^0$, no higher cohomology intervening);
\item $H^i(X,L^{\otimes n})=0$ for all $i>0$ and $n>0$;
\item the global sections are discrete over $A^{\triangleright}(\ast)$:
$R\Gamma(X,L^{\otimes n})\in D_{\mathrm{disc}}(A^{\triangleright}(\ast))\subseteq D(A)$.
\end{enumerate}
Then
\begin{equation}
\label{eq:21:xalg}
X^{\mathrm{alg}}\ :=\ \Proj\ \bigoplus_{n\ge0} R\Gamma(X,L^{\otimes n})\ \in\
\mathrm{DerSch}_{A^{\triangleright}(\ast)}
\end{equation}
is a quasicompact quasiseparated proper scheme over $A^{\triangleright}(\ast)$,
with a proper map $f\colon X\to X^{\mathrm{alg}}_A$. If $L$ is ample enough that
$f$ is an isomorphism, then $X\xrightarrow{\sim}X^{\mathrm{alg}}_A$ is the
analytification of $X^{\mathrm{alg}}$.
\end{proposition}

\begin{remark}[$f_*\mathcal O=\mathcal O$]
\label{rmk:21:pushforward-unit}
The crucial observation is that $f_*(\mathbf 1)=\mathbf 1$, i.e.\
$f_*\mathcal O_X=\mathcal O_{X^{\mathrm{alg}}_A}$. This may be checked after
pullback to an affine chart $U^{\mathrm{alg}}_A=D^{+}_1(\xi)\subseteq
X^{\mathrm{alg}}_A$ (the non-vanishing locus of a section
$\xi\in H^0(X,L^{\otimes n})$), and there on global sections. The relevant global
sections are a filtered colimit along multiplication by $\xi$,
\begin{equation}
\label{eq:21:section-colimit}
\varinjlim_{k,\,\times\xi} R\Gamma\bigl(X,L^{\otimes kn}\bigr)\ =\
\varinjlim_{k,\,\times\xi} R\Gamma\bigl(X^{\mathrm{alg}},L^{\otimes kn}\bigr),
\end{equation}
i.e.\ the localisation inverting $\xi$; since $X^{\mathrm{alg}}$ was built from
exactly these graded pieces, the two sides agree, giving $f_*\mathcal O=\mathcal O$.
\end{remark}

\begin{remark}[From $f_*\mathcal O=\mathcal O$ to an isomorphism]
\label{rmk:21:iso-criterion}
Because $f$ is proper and $f_*(\mathbf 1)=\mathbf 1$, the map $f$ is surjective:
$f_*(\mathbf 1)=\mathbf 1$ is preserved by every base change (proper base change),
so after pullback to any affine the target ring embeds into the pushed-forward
sections, forcing $f$ surjective. To upgrade surjectivity to an isomorphism it
suffices that the diagonal $\Delta_f$ be an isomorphism (a surjection of sheaves
with iso diagonal is iso). The identity $f_*(\mathbf 1)=\mathbf 1$ also shows,
via the projection formula, that the pullback
$f^{*}\colon D(X^{\mathrm{alg}}_A)\hookrightarrow D(X)$ is fully faithful (the
unit of $f^{*}\dashv f_{*}$ is an equivalence, $f_*f^{*}(-)=(-)\otimes f_*\mathbf 1
=(-)$), and likewise after any base change. Hence it is enough to show that the
structure sheaf of the diagonal $\Delta_X$---a coherent sheaf on
$X^{\mathrm{alg}}_A\times X^{\mathrm{alg}}_A$, supported on a Zariski-closed
immersion---lies in the essential image of $f^{*}\times f^{*}$; then it is pulled
back from $X^{\mathrm{alg}}\times X^{\mathrm{alg}}$ and $\Delta_f$ is an
isomorphism.
\end{remark}

\begin{remark}[The case of $E_q$]
\label{rmk:21:tate-algebraisation}
For the Tate curve one takes for $L$ the ample line bundle attached to the zero
(unit) section. All global sections $R\Gamma(E_q,L^{\otimes n})$ are then
computed inductively---$R\Gamma(E_q,\mathcal O)$ directly, and each further power
adds copies of the base ring from the zero section---and the required diagonal
condition is verified using the group structure: the diagonal is translated onto
the zero section, and the zero section is the vanishing locus one already
controls. This produces the ample bundle and the isomorphism
$E_{q,A}\xrightarrow{\sim}E_q^{\mathrm{alg}}$, proving
\cref{thm:21:tate-algebraic}.
\end{remark}

\begin{remark}[Properness here is not Grothendieck's]
\label{rmk:21:properness-comparison}
In the closing discussion Scholze stressed that the notion of proper used above
(\cref{def:21:proper}) is not literally Grothendieck's finite-type + separated +
universally closed. In particular, for schemes embedded into analytic stacks
\emph{every} map is proper in this sense; so the section ring
$\bigoplus_n R\Gamma(X,L^{\otimes n})$ may be far from finitely generated (its
degree-zero part alone can be an arbitrary ring). When $X$ genuinely comes from a
more classical space---for instance a complex-analytic space, which embeds into
analytic stacks---the analytic notion of properness specialises to the usual one
(it implies quasicompactness), and algebraisation can succeed only when $X$ is
proper in the universally-closed sense. The GAGA-type isomorphisms in this world
hold for these ``funny'' analytic stacks, of which schemes and analytic spaces
are full subcategories.
\end{remark}
